\chapter{Correctness Arguments}
\label{ch:correctness}

All six algorithms operate on the state-expanded graph $G^{*}$ under the variable purchase
model: the driver chooses exactly how much fuel $b \in [0, C-f]$ to buy at each station, and
the state $(v,f)$ tracks both position and remaining fuel.

\section{Dijkstra-Based Approach (Baseline 1)}

Correctness follows from two properties of $G^{*}$:
\begin{enumerate}
    \item \textbf{Non-negative edge costs.}
          All travel edges have cost~0 and all refuel edges have cost $b \cdot p(v) \geq 0$,
          satisfying Dijkstra's prerequisite.
    \item \textbf{Optimal substructure.}
          If the optimal path passes through state $(v,f)$, the sub-path to $(v,f)$ must also
          be optimal, since fuel costs decompose additively over refuel edges.
\end{enumerate}
Because Dijkstra exhausts all states $(v,f)$ in non-decreasing cost order, it is guaranteed
to find the globally minimum purchase cost for any $f_0$.

\section{Greedy Strategy: Counterexample (Baseline 2)}

The greedy strategy is not optimal in general.
On the running example with $\eta=10$\,km/L, $C=40$\,L, $f_0=28$\,L:

\begin{table}[ht]
\centering
\caption{Greedy counterexample.}
\label{tab:greedy_counter}
\begin{tabular}{llll}
\toprule
Strategy & Buy at A & Buy at E & Total \\
\midrule
Greedy  & 12\,L@36 $= 432$\,THB & 18\,L@24 $= 432$\,THB & 864\,THB ($+300\%$) \\
Optimal & 0\,L & 9\,L@24 $= 216$\,THB & 216\,THB \\
\bottomrule
\end{tabular}
\end{table}

The failure is structural: greedy cannot see that $f_0=28$\,L already suffices to coast to
station E without any purchase at A.

\section{Dynamic Programming Correctness (Baseline 3)}

DP correctness follows from the principle of optimality applied to the variable-purchase
state space.
The recurrence
\[
    \mathrm{DP}[w][f_w] = \min_{b}\!\bigl(\mathrm{DP}[v][f] + b \cdot p(v)\bigr)
\]
considers every purchase amount $b \in \{0,\varepsilon,\ldots,C-f\}$ at every predecessor
state $(v,f)$ and propagates the minimum cost forward.
Initializing with $\mathrm{DP}[s][f_0]=0$ for arbitrary $f_0$ correctly captures a partial
starting tank.
Taking $\min_{f}\mathrm{DP}[t][f]$ over all arrival fuel levels yields the true global
optimum.

\section{Standard A* Correctness (Baseline 4)}

A* uses $h(v) = d_{\mathrm{geo}}(v,t)/\eta \cdot p_{\min}$.

\paragraph{Admissibility.}
For any state $(v,f)$:
\[
    h(v) = \frac{d_{\mathrm{geo}}(v,t)}{\eta} \cdot p_{\min}
         \leq \frac{d_{\mathrm{road}}(v,t)}{\eta} \cdot p_{\min}
         \leq h^{*}(v,f),
\]
since straight-line distance $\leq$ road distance, and $p_{\min}$ is the global lower bound
on any fuel price paid.

\paragraph{Consistency.}
For any edge $(v,f) \to (w,f')$ in $G^{*}$:
\[
    h(v) \leq C(v,w) + h(w)
\]
by the triangle inequality on geodesic distances, so the heuristic satisfies the consistency
condition.

Both properties guarantee that A* returns an optimal solution.
However, because $h$ ignores $f$, states $(v,f_1)$ and $(v,f_2)$ with $f_1 \neq f_2$ receive
identical heuristic values, limiting pruning effectiveness.

\section{RF-A* Correctness (Baseline 5)}

RF-A* uses the same heuristic $h_{\mathrm{RF}}(v)$ as standard A* and inherits the same
admissibility and consistency arguments.
Optimality of each segment is guaranteed by the underlying A* call; the segment-by-segment
construction preserves overall optimality because refueling costs compose additively.

\section{PF-A* Correctness}
\label{sec:pfa_correctness}

PF-A* replaces $h_{\mathrm{RF}}$ with the fuel-level-aware heuristic
\[
    h_{\mathrm{PF}}(v, f)
    = \max\!\left(\frac{d_{\mathrm{geo}}(v, t)}{\eta} - f,\; 0\right)
      \cdot p_{\min}^{\mathrm{reach}}(v, f).
\]

\subsection{Admissibility}

\begin{lemma}[Fuel-need underestimate]
    $\max(d_{\mathrm{geo}}(v,t)/\eta - f,\; 0) \leq$ actual additional fuel required from
    state $(v,f)$ on any feasible path.
\end{lemma}
\begin{proof}
    Let $\ell^{*}$ be the road distance of the optimal remaining path.
    Since $d_{\mathrm{geo}}(v,t) \leq \ell^{*}$, the actual fuel needed beyond $f$ satisfies
    $\ell^{*}/\eta - f \geq d_{\mathrm{geo}}(v,t)/\eta - f$.
    Taking $\max(\cdot, 0)$ on both sides preserves the inequality.
\end{proof}

\begin{lemma}[Price underestimate]
    $p_{\min}^{\mathrm{reach}}(v,f) \leq$ any fuel price paid on the optimal remaining path
    from state $(v,f)$.
\end{lemma}
\begin{proof}
    Every station visited on the optimal path from $(v,f)$ must be reachable from $v$
    given at least $f$ litres of fuel.
    By definition, $p_{\min}^{\mathrm{reach}}(v,f)$ is the minimum price over all such
    stations, so it is a lower bound on every price paid.
\end{proof}

\begin{theorem}[Admissibility of $h_{\mathrm{PF}}$]
    $h_{\mathrm{PF}}(v,f) \leq h^{*}(v,f)$ for all states $(v,f)$.
\end{theorem}
\begin{proof}
    Let $\Delta f = \max(d_{\mathrm{geo}}(v,t)/\eta - f,\; 0)$ and let $p^{*}$ be the price
    at which the additional fuel is actually purchased on the optimal path.
    By Lemma~1, the actual additional fuel needed $\geq \Delta f$.
    By Lemma~2, $p^{*} \geq p_{\min}^{\mathrm{reach}}(v,f)$.
    Therefore $h^{*}(v,f) \geq \Delta f \cdot p^{*} \geq \Delta f \cdot p_{\min}^{\mathrm{reach}}(v,f) = h_{\mathrm{PF}}(v,f)$.
\end{proof}

\subsection{Consistency}

\begin{theorem}[Consistency of $h_{\mathrm{PF}}$]
    For every transition $(v,f) \to (w,f')$ in $G^{*}$ with transition cost $c$:
    \[
        h_{\mathrm{PF}}(v,f) \leq c + h_{\mathrm{PF}}(w,f').
    \]
\end{theorem}
\begin{proof}
    We consider the two types of transitions separately.

    \textbf{Travel edge} $(v,f) \to (w,\; f - d(v,w)/\eta)$, cost $c = 0$.
    Let $\delta = d(v,w)/\eta$ (fuel consumed).
    Then $f' = f - \delta$.
    \begin{align*}
        h_{\mathrm{PF}}(w, f')
        &= \max\!\left(\frac{d_{\mathrm{geo}}(w,t)}{\eta} - f',\; 0\right)
           \cdot p_{\min}^{\mathrm{reach}}(w, f') \\
        &\geq \max\!\left(\frac{d_{\mathrm{geo}}(w,t)}{\eta} - f',\; 0\right)
           \cdot p_{\min}^{\mathrm{reach}}(v, f)
           \quad\text{(reachable set from $w$ may be smaller)}.
    \end{align*}
    By the triangle inequality on geodesic distances,
    $d_{\mathrm{geo}}(v,t) \leq d(v,w) + d_{\mathrm{geo}}(w,t)$, so
    $d_{\mathrm{geo}}(v,t)/\eta - f \leq d_{\mathrm{geo}}(w,t)/\eta - f'$.
    Hence $h_{\mathrm{PF}}(v,f) \leq h_{\mathrm{PF}}(w,f') = 0 + h_{\mathrm{PF}}(w,f')$.

    \textbf{Refuel edge} $(v,f) \to (v,\; f+b)$, cost $c = b \cdot p(v)$, $b > 0$.
    Then $f' = f + b$.
    \begin{align*}
        c + h_{\mathrm{PF}}(v,f')
        &= b \cdot p(v) + \max\!\left(\frac{d_{\mathrm{geo}}(v,t)}{\eta} - f - b,\; 0\right)
           \cdot p_{\min}^{\mathrm{reach}}(v, f').
    \end{align*}
    Let $D = d_{\mathrm{geo}}(v,t)/\eta$.
    \textit{Case 1}: $f \geq D$.
    Then $h_{\mathrm{PF}}(v,f) = 0 \leq c + h_{\mathrm{PF}}(v,f')$.
    \textit{Case 2}: $f < D \leq f + b$.
    Then $h_{\mathrm{PF}}(v,f) = (D-f)\cdot p_{\min}^{\mathrm{reach}}(v,f)$ and
    $h_{\mathrm{PF}}(v,f') = 0$, so we need $(D-f)\cdot p_{\min}^{\mathrm{reach}}(v,f) \leq b \cdot p(v)$.
    Since $D - f \leq b$ and $p_{\min}^{\mathrm{reach}}(v,f) \leq p(v)$,
    this holds.
    \textit{Case 3}: $f + b < D$.
    Then both sides are positive.
    We need $(D-f)\cdot p_{\min}^{\mathrm{reach}}(v,f) \leq b\cdot p(v) + (D-f-b)\cdot p_{\min}^{\mathrm{reach}}(v,f')$.
    Since $p_{\min}^{\mathrm{reach}}(v,f') \geq p_{\min}^{\mathrm{reach}}(v,f)$ (buying more
    fuel does not extend range to cheaper stations) and $p(v) \geq p_{\min}^{\mathrm{reach}}(v,f)$,
    the right-hand side dominates.
\end{proof}

\subsection{Relationship to RF-A*}

It is tempting to claim $h_{\mathrm{PF}}(v,f) \leq h_{\mathrm{RF}}(v)$ for all $(v,f)$, but
this does \emph{not} hold universally.
The two heuristics differ along two axes that can push in opposite directions: the
fuel-need term $\max(d_{\mathrm{geo}}(v,t)/\eta - f,\; 0)$ in $h_{\mathrm{PF}}$ is always
no larger than the corresponding term $d_{\mathrm{geo}}(v,t)/\eta$ in $h_{\mathrm{RF}}$,
but the price term $p_{\min}^{\mathrm{reach}}(v,f)$ can exceed the global $p_{\min}$.
A universal pointwise dominance between $h_{\mathrm{PF}}$ and $h_{\mathrm{RF}}$ would
therefore require an additional assumption relating the price ratio to the fuel ratio, and we
do not claim such an assumption holds in general.

Instead, we make a weaker but defensible claim.
$h_{\mathrm{PF}}$ is \emph{state-aware}: it uses the current fuel level $f$ and the set of
stations actually reachable from $(v,f)$, whereas $h_{\mathrm{RF}}$ depends only on $v$.
This additional information allows $h_{\mathrm{PF}}$ to distinguish states that share the
same $v$ but differ in $f$ or in reachable-station set, which RF-A* cannot do.
On networks with non-trivial price variance and sufficient station density, which is the regime of
interest for real Thai corridors, this state-awareness empirically reduces node expansions,
as reported in Chapters~\ref{ch:complexity} and~\ref{ch:results}.

\begin{proposition}[State-awareness of $h_{\mathrm{PF}}$]
    For any fixed node $v$, the map $f \mapsto h_{\mathrm{PF}}(v,f)$ is non-increasing in
    $f$, whereas $h_{\mathrm{RF}}(v)$ is constant in $f$.
\end{proposition}
\begin{proof}
    Fix $v$ and let $f_1 \leq f_2$.
    The fuel-need term satisfies
    $\max(d_{\mathrm{geo}}(v,t)/\eta - f_1,\; 0)
     \geq \max(d_{\mathrm{geo}}(v,t)/\eta - f_2,\; 0)$.
    The reachable-station set from $(v,f_2)$ contains the reachable-station set from
    $(v,f_1)$, so
    $p_{\min}^{\mathrm{reach}}(v,f_2) \leq p_{\min}^{\mathrm{reach}}(v,f_1)$.
    Both factors are non-negative and non-increasing in $f$, so their product is
    non-increasing in $f$.
    The constancy of $h_{\mathrm{RF}}(v)$ is immediate from its definition.
\end{proof}

Because $h_{\mathrm{PF}}$ is both admissible and consistent (Theorems above), PF-A*
inherits the optimality guarantee of RF-A*.
The empirical benefit of PF-A* over RF-A* in node expansions is evaluated directly in
Experiments~5 and~6.
